
\section{Related Work}
\textbf{Temporal Graph Learning}
Temporal graph learning has been widely studied for dynamic link prediction. 
TGAT combines temporal neighborhood attention with functional time encoding \cite{xu2020tgat}. 
TGN provides a memory-based continuous time framework and remains a strong baseline \cite{rossi2020tgn}. 
Recent methods further improve efficiency or long-history modeling, including GraphMixer, DyGFormer, and RepeatMixer \cite{cong2023graphmixer,yu2023dygformer,zou2024repeatmixer}, 
while newer studies revisit temporal position encoding, target-aware matching, and interaction heterogeneity \cite{cheng2025scalability,tieu2025learnable,yi2026future,yu2026tami}. 
Heterogeneous graph models such as HAN and HGT capture typed nodes and relations, and temporal variants such as STHN extend this idea to dynamic typed graphs \cite{wang2019heterogeneous,hu2020hgt,li2023sthn,yue2025heterogeneous}. 
Benchmark efforts, including TGB, TGB 2.0, and TGB-Seq, provide realistic protocols for temporal, multi-relational, and future-link prediction \cite{huang2023tgb,gastinger2024tgb2,yi2025tgbseq}. 
These studies motivate our chronological evaluation and strong temporal baselines. However, many methods still treat relation types as static edge features. 
DHC follows a strict chronological protocol, but instead focuses on relation-aware dynamic temporal context by explicitly maintaining relation memory in heterogeneous streams for link scoring.
\rev{The closest architectural baseline, STHN, simplifies continuous-time heterogeneous link prediction by sampling a fixed number of recent interactions and encoding them with a compact heterogeneous temporal network. DHC differs at the state and query levels: it assigns every relation type a recurrent state updated from chronological messages, uses that state to gate endpoint messages and neighbor-edge attributes, and fuses the resulting query-specific temporal subgraph representation with an independently maintained TGN event memory. TAMI instead targets interaction heterogeneity through target-aware matching; it does not instantiate DHC's recurrent per-relation state. Thus, DHC's intended novelty is dynamic relation semantics and their use throughout message generation, local context encoding, and scoring, rather than temporal sampling alone.}

\textbf{Recent Link Prediction on Academic Networks}
Academic networks have become representative testbeds for temporal heterogeneous link prediction. 
Recent studies model papers, authors, venues, datasets, patents, assignees, and technology classes as typed entities, and cast prediction tasks as forecasting future citation, collaboration, publication--dataset, or patent impacts. 
For instance, publication--dataset link prediction on sparse scholarly knowledge graphs highlights the importance of inductive generalization and data discoverability \cite{irrera2025heterogeneous}.
Citation-oriented studies further exploit dynamic academic graphs, clustering structures, centrality signals, and bias-aware graph representations to estimate future scientific impact \cite{zhang2025cdchgn,lu2026bacite}.
Similarly, patent-related studies formulate citation and technology-impact prediction as dynamic learning problems over technical indicators and evolving citation patterns \cite{seol2026patentcitation}.
Beyond these settings, recent work has extended the empirical scope to topology-aware link prediction on citation benchmarks, heterogeneous link prediction in biomedical networks, and patent-based technology-opportunity discovery \cite{hu2025biomedicalhgnn,wu2026identifying,liang2026tagnn}
Other recent work broadens the evidence base through topology-focused link prediction on citation benchmarks, heterogeneous link prediction in biomedical networks, and patent-based technology-opportunity discovery \cite{hu2025biomedicalhgnn,wu2026identifying,liang2026tagnn}. 
Collectively, these studies motivate the design of DHC, particularly its emphasis on preserving typed contextual edges, relation states, and strictly chronological evaluation.
